%%%%%%%%%%%%%%%%%%%%%%%%%%%%%%%%%%%%%%%%%
% "ModernCV" CV y Carta de Presentación
% Plantilla en LaTeX
% Versión 1.11 (19/6/14)
%
% Esta plantilla ha sido descargada de:
% http://www.LaTeXTemplates.com
%
% Autor original:
% Xavier Danaux (xdanaux@gmail.com)
%
% Licencia:
% CC BY-NC-SA 3.0 (http://creativecommons.org/licenses/by-nc-sa/3.0/)
%
% Nota importante:
% Esta plantilla requiere que los archivos moderncv.cls y .sty estén en el 
% mismo directorio que este archivo .tex. Estos archivos proporcionan el estilo 
% y los temas utilizados para estructurar el documento.
%
%%%%%%%%%%%%%%%%%%%%%%%%%%%%%%%%%%%%%%%%%

%----------------------------------------------------------------------------------------
%	PAQUETES Y CONFIGURACIONES DEL DOCUMENTO
%----------------------------------------------------------------------------------------

\documentclass[12pt,a4paper,sans]{moderncv} % Tamaños de fuente: 10, 11 o 12; tamaños de papel: a4paper, letterpaper, a5paper, legalpaper, executivepaper o landscape; familias tipográficas: sans o roman.

\usepackage[utf8]{inputenc}

\usepackage[export]{adjustbox}

\usepackage{multicol}

\moderncvstyle{casual} % Tema del CV - opciones: 'casual' (por defecto), 'classic', 'oldstyle' y 'banking'.
\moderncvcolor{black} % Color del CV - opciones: 'blue' (por defecto), 'orange', 'green', 'red', 'purple', 'grey' y 'black'.

\usepackage{color}

\usepackage{lipsum} % Para insertar texto de relleno 'Lorem ipsum'.

\usepackage[scale=0.85]{geometry} % Reduce los márgenes del documento.
\setlength{\hintscolumnwidth}{4cm} % Ajusta el ancho de la columna de fechas.
%\setlength{\makecvtitlenamewidth}{10cm} % Para el estilo 'classic', ajustar el ancho asignado al nombre.

%----------------------------------------------------------------------------------------
%	NOMBRE Y DATOS DE CONTACTO
%----------------------------------------------------------------------------------------

\firstname{Juan Pablo} % Tu nombre.
\familyname{Liévano Karim} % Tu apellido.

%----------------------------------------------------------------------------------------

\begin{document}

\Huge{Juan Pablo Liévano Karim}\normalsize
\hrulefill
%----------------------------------------------------------------------------------------
%	DATOS DE CONTACTO
%----------------------------------------------------------------------------------------

\href{mailto:lievanojuanpablo@gmail.com}{lievanojuanpablo@gmail.com}


%----------------------------------------------------------------------------------------
%	EDUCACIÓN
%----------------------------------------------------------------------------------------

\section{Educación}
\cvitem{2018 -- 2020}{\textbf{Maestría en Matemáticas, IMPA.} \textit{Río de Janeiro, Brasil}.}
\cvitem{2014 -- 2018}{\textbf{Pregrado en Matemáticas, Universidad de los Andes.} \textit{Bogotá, Colombia} \newline - Opción en Literatura. \newline - Opción en Historia de la Música.}
\cvitem{2025}{\textbf{BlueDot Impact} (San Francisco). Dos certificaciones (q.v.\ bluedot.org): \newline - \textit{AGI Strategy}. \newline - \textit{Technical AI Safety}.}

%----------------------------------------------------------------------------------------
%	EXPERIENCIA
%----------------------------------------------------------------------------------------

\section{Experiencia}
\cvitem{}{\small\textit{Experiencias seleccionadas, organizadas según su relevancia profesional.}}
\cvitem{2024, 2025}{\textbf{Investigador en CHAI, UC Berkeley} (dos veces). \textit{El laboratorio de IA segura de Berkeley, dirigido por Stuart Russell.} \newline - 2025: escribí \href{https://openreview.net/forum?id=XpJgFwSJFM}{un artículo} (asesorado por \href{https://camallen.net}{Cameron Allen}) con experimentos que muestran que podemos decodificar e interpretar las activaciones neuronales dentro de redes neuronales recurrentes en ciertos entornos. \newline - 2024: programa ``Asking for Help'' de alineación de IA (con \href{https://bplaut.github.io}{Benjamin Plaut}), que derivó en el artículo coautorado con Stuart Russell.}
\cvitem{2025, 2026}{\textbf{Docente, Universidad de los Andes}. \textit{Facultad de Economía.} Profesor de dos cursos avanzados: \newline - ``Machine Learning para Business Intelligence''. \newline - ``Deep Learning y aplicaciones avanzadas de Machine Learning''.}
\cvitem{2026}{\textbf{Facilitador de curso, BlueDot Impact}. Instructor del curso Technical AI Safety.}
\cvitem{2023 -- 2025}{\textbf{Investigador en ML y desarrollador Python, Quantil}. Firma colombiana líder en matemáticas aplicadas e IA, con clientes como Ecopetrol, Banco W y el Estado colombiano.}
\cvitem{2018, 2021, 2022}{\textbf{Editor y traductor de currículo, \textit{Illustrative Mathematics}}. Organización líder en educación matemática, cuyos currículos basados en evidencia se usan internacionalmente.}
\cvitem{2016, 2017}{\textbf{Profesor complementario, Universidad de los Andes}. \textit{Departamento de Matemáticas.} Dos cursos introductorios: \newline - Álgebra Lineal. \newline - Cálculo Integral.}
\cvitem{2023}{\textbf{Asistente de investigación, UCLA}. Análisis cualitativo de datos y edición conceptual para un proyecto de bienestar social con madres adolescentes en Cartagena y Medellín.}
\cvitem{2022 -- actual}{\textbf{Consultor de datos, Joyería Liévano}. Modelos predictivos y análisis de clientes.}

%----------------------------------------------------------------------------------------
%	RECONOCIMIENTOS, PROYECTOS Y CONFERENCIAS POR INVITACIÓN
%----------------------------------------------------------------------------------------

\section{Reconocimientos, proyectos y conferencias por invitación}

\cvitem{2026}{\textbf{Jurado de AI Safety Colombia}. Elegido para formar parte del panel de jurados del \href{https://aisafetycolombia.org/en/hackathon/}{AI Safety Colombia Hackathon 2026} organizado por \href{https://apartresearch.com}{Apart Research}.}

\cvitem{2025}{\textbf{Primer autor en NeurIPS}. \textit{Una de las conferencias más importantes del mundo en inteligencia artificial.} Presenté el artículo \textit{``Echo of Bayes: Learned Memory Functions Can Recover Belief States''} (con UC Berkeley y Brown University). Disponible en \href{https://openreview.net/forum?id=XpJgFwSJFM}{OpenReview}.}

\cvitem{2025}{\textbf{Coautor de un artículo con \href{https://people.eecs.berkeley.edu/~russell/}{Stuart Russell}}. (Reconocido cientifico y autor de «Artificial Intelligence: A Modern Approach», el libro de texto de IA más usado del mundo). Articulo: \textit{``Safe Learning Under Irreversible Dynamics via Asking for Help''} \href{https://arxiv.org/abs/2502.14043}{(arXiv:2502.14043)}.}

\cvitem{2026}{\textbf{Artículo en preparación} sobre sistemas de recomendación, con Sana Panday (UC Berkeley): diseñar algoritmos de recomendación que incorporen incentivos de construcción de paz, en lugar de optimizar únicamente la captación de atención (\textit{engagement}).}

\cvitem{2018}{\textbf{Beca completa con estipendio (CNPq)} para la maestría en IMPA.}

\cvitem{2024, 2025, 2026}{\textbf{Participante invitado al taller anual de CHAI} en Asilomar, California. Retiro por invitación que reúne cada año a unos 200 investigadores líderes en IA y seguridad de IA.}


%----------------------------------------------------------------------------------------
%	HABILIDADES
%----------------------------------------------------------------------------------------

\section{Habilidades}
\cvitem{}{- Python (Pandas para ciencia de datos; JAX/PyTorch para aprendizaje por refuerzo).}{}
\cvitem{}{- Amplio conocimiento en estadística, probabilidad, álgebra lineal y lógica formal.}{}
\cvitem{}{- Pedagogía matemática: diseño de currículo (Illustrative Mathematics) y docencia universitaria (Uniandes).}{}
\cvitem{}{- Infraestructura y datos: Azure (OCR, Cognitive Search, Cosmos DB), Databricks, \LaTeX.}{}
\cvitem{}{- Experiencia en diseño y desarrollo web.}

%----------------------------------------------------------------------------------------
%	IDIOMAS
%----------------------------------------------------------------------------------------

\section{Idiomas}

\cvitem{}{Español, \textit{Nativo}.}{}
\cvitem{}{Inglés, \textit{Muy avanzado}.}{}
\cvitem{}{Portugués, \textit{Avanzado}.}{}

%----------------------------------------------------------------------------------------

\end{document}
